\section*{Chapter \ref{ch:iv:linear-algebra-and-calculus} --- Linear Algebra and Calculus}

\begin{solution}{iq:iv:linear-algebra-and-calculus:1}
The matrix is $I + \mathbf 1\mathbf 1^\top$ in two dimensions: eigenvalue 3 with eigenvector $(1, 1)$ and 1 with $(1,
-1)$. Trace 4 and determinant 3 confirm.
\iqlookfor{recognising the structure and checking with trace and determinant.}
\end{solution}

\begin{solution}{iq:iv:linear-algebra-and-calculus:2}
No. By \cref{prop:iv:linear-algebra-and-calculus:third}, with $\rho_{12} = \rho_{13} = 0.9$ the third correlation
must lie in $[0.81 - 0.19, 0.81 + 0.19] = [0.62, 1]$. The determinant is $1 - 3 \times 0.81 + 2 \times 0.9 \times 0.9
\times (-0.9) = -2.888 < 0$, and the eigenvalues are $-0.8$, 1.9 and 1.9. Intuitively, assets 2 and 3 cannot both
move with asset 1 and against each other.
\iqlookfor{a feasibility check by the bound or the determinant, with the intuition.}
\end{solution}

\begin{solution}{iq:iv:linear-algebra-and-calculus:3}
Completing the square, $\E[e^Z] = e^{1/2} \approx 1.649$ (the lognormal mean). $\E[Z^4] = 3$ (the fourth moment of a
normal, used as the benchmark kurtosis).
\iqlookfor{the moment generating function and the normal kurtosis.}
\end{solution}

\begin{solution}{iq:iv:linear-algebra-and-calculus:4}
Differentiate $\sum x^k = 1/(1-x)$: $\sum kx^k = x/(1-x)^2$, which at $x = \tfrac12$ is 2. (It is also the mean of
a geometric waiting time.)
\iqlookfor{differentiating a known series.}
\end{solution}

\begin{solution}{iq:iv:linear-algebra-and-calculus:5}
$0.48 \pm \sqrt{(1 - 0.64)(1 - 0.36)} = 0.48 \pm 0.48$: any value from 0 to 0.96.
\iqlookfor{the determinant condition applied with numbers.}
\end{solution}

\begin{solution}{iq:iv:linear-algebra-and-calculus:6}
Eigenvalue $1 + 9\rho$ (eigenvector $\mathbf 1$) and $1 - \rho$ with multiplicity 9. Non-negativity needs $\rho \ge
-1/9$: ten assets can be at most about $-0.11$ correlated with each other on average, because their equally weighted sum
cannot have negative variance.
\iqlookfor{the equicorrelation spectrum and the portfolio interpretation of the bound.}
\end{solution}

\begin{solution}{iq:iv:linear-algebra-and-calculus:7}
$H$ is a projection onto a five-dimensional space: eigenvalues 1 (five times) and 0 (245 times), trace 5. The average
leverage is $5/250 = 0.02$; observations with leverage several times that (extreme regressor values) pull the fit
towards themselves and deserve a look, as does any regression whose $p/n$ is not small.
\iqlookfor{the projection's spectrum and the meaning of leverage.}
\end{solution}

\begin{solution}{iq:iv:linear-algebra-and-calculus:8}
$w_1 = (0.01 - 0.3 \times 0.02)/(0.04 + 0.01 - 2 \times 0.3 \times 0.02) = 0.004/0.038 \approx 0.105$ in the
volatile asset, 0.895 in the other, for a volatility of about 9.8\%, lower than the less volatile asset's 10\% because the
correlation is below $\sigma_2/\sigma_1 = 0.5$.
\iqlookfor{the two-asset formula and the diversification condition.}
\end{solution}

\begin{solution}{iq:iv:linear-algebra-and-calculus:9}
$\E[Z^+] = \int_0^\infty z\varphi(z)\,dz = \varphi(0) = 1/\sqrt{2\pi} \approx 0.399$. It gives the at-the-money
approximation of an option: a forward at-the-money call is worth about $0.4\,S\sigma\sqrt T$
(\cref{ch:iv:options-and-derivatives}).
\iqlookfor{the integral by the density's derivative, and the link to option prices.}
\end{solution}

\begin{solution}{iq:iv:linear-algebra-and-calculus:10}
By symmetry the answer keeps the pattern of signs with equal magnitudes; alternating projections (or eigenvalue
clipping, which coincides here) give $\rho_{12} = \rho_{13} = 0.5$ and $\rho_{23} = -0.5$, with eigenvalues 0, 1.5, 1.5:
every entry moves by 0.4, a Frobenius distance of $\sqrt{6 \times 0.16} \approx 0.98$. For a thousand assets use the
alternating-projections algorithm with Dykstra's correction (each step one eigendecomposition, $O(n^3)$) or a Newton
method on the dual (Higham's paper and its successors), and report which entries moved.
\iqlookfor{the structure of the answer, its distance, and the algorithm at scale.}
\end{solution}

\begin{solution}{iq:iv:linear-algebra-and-calculus:11}
$\|v\|^2 = 9$: eigenvalues 10 (along $v$) and 1, 1. The inverse is $I - vv^\top/10$ by Sherman--Morrison. A one-factor
covariance $\sigma^2 I + \beta\beta^\top$ is a scaled version, so its inverse (needed for minimum-variance weights or a
Mahalanobis distance) costs $O(n)$, not $O(n^3)$.
\iqlookfor{the spectrum of a rank-one update and the computational payoff.}
\end{solution}

\begin{solution}{iq:iv:linear-algebra-and-calculus:12}
Maximise $\mu^\top w$ subject to $w^\top\Sigma w = 0.01$: $w \propto \Sigma^{-1}\mu = (0.05/0.04, 0.03/0.01) = (1.25,
3)$, whose volatility is $\sqrt{1.25^2 \times 0.04 + 3^2 \times 0.01} \approx 0.39$; scale by $0.1/0.39$ to $w \approx
(0.320, 0.768)$. Expected return about 3.91\%, Sharpe ratio about 0.39, higher than either asset's 0.25 and 0.30.
\iqlookfor{the Lagrangian solution $\Sigma^{-1}\mu$ and the scaling.}
\end{solution}

\begin{solution}{iq:iv:linear-algebra-and-calculus:13}
The Lagrange remainder is $e^\xi x^3/6$ for some $\xi \in (0, x)$, at most $e^{0.1} \times 0.001/6 \approx 0.000184$. The
actual error is $e^{0.1} - 1.105 \approx 0.000171$, inside the bound.
\iqlookfor{the remainder term, bounded honestly.}
\end{solution}
